% shared/notation.tex — central registry of mathematical notation.
% All math symbols used in en/ja/zh chapters MUST be defined here.
% Gate 3 (notation registration) greps chapter sources against this file.
%
% Naming conventions:
%   * scalars       lowercase     (r, b, t)
%   * vectors       \mathbf{...}  (\mathbf{e})
%   * sets/spaces   \mathcal{...} (\mathcal{N}, \mathcal{T})
%   * functions     \mathrm{...}  (\mathrm{ch}, \mathrm{pa})

% ============================== sets / spaces
\newcommand{\Nodes}{\mathcal{N}}
\newcommand{\Tree}{\mathcal{T}}
\newcommand{\Reals}{\mathbb{R}}
\newcommand{\Naturals}{\mathbb{N}}

% ============================== individual nodes
\newcommand{\node}[1]{n_{#1}}
\newcommand{\children}[1]{\mathrm{ch}(#1)}
\newcommand{\parentof}[1]{\mathrm{pa}(#1)}
\newcommand{\depth}[1]{\mathrm{depth}(#1)}

% ============================== node state
\newcommand{\statefn}[1]{s(#1)}
% State labels are usable in both math and text mode (the body sometimes
% references them inline outside $...$). \ensuremath wraps the math content.
\newcommand{\StateOpen}{\ensuremath{\mathrm{open}}}
\newcommand{\StateEval}{\ensuremath{\mathrm{eval}}}
\newcommand{\StateDone}{\ensuremath{\mathrm{done}}}
\newcommand{\StateFailed}{\ensuremath{\mathrm{failed}}}

% ============================== rewards and evaluation axes
\newcommand{\nodescore}[1]{r(#1)}            % scalar reward
\newcommand{\axes}[1]{\mathbf{e}(#1)}        % per-axis evaluation vector
\newcommand{\aggregate}{\phi}                % \phi: R^k -> R
\newcommand{\axesdim}{k}                     % number of axes

% ============================== exploration utility / policy
\newcommand{\utility}[1]{U(#1)}
\newcommand{\rewardhat}[1]{\hat r(#1)}
\newcommand{\novelty}[1]{\mathrm{novelty}(#1)}
\newcommand{\policy}{\pi}
\newcommand{\softmaxbeta}{\beta}
\newcommand{\utilcoef}[1]{\alpha_{#1}}

% ============================== budget
\newcommand{\bbranch}{b}
\newcommand{\Tmax}{T_{\max}}
\newcommand{\Cmax}{C_{\max}}

% ============================== rubric
\newcommand{\Rubric}{R}
\newcommand{\rubricitem}[1]{c_{#1}}
\newcommand{\rubricweight}[1]{w_{#1}}
\newcommand{\judgefn}[1]{\mathrm{judge}_{#1}}

% ============================== paper score
\newcommand{\paper}{P}
\newcommand{\paperscore}[1]{\mathrm{score}(#1)}
\newcommand{\refine}{\mathrm{Refine}}

% ============================== checkpoint / lineage
\newcommand{\ckpt}{\mathrm{ckpt}}
\newcommand{\lineage}[1]{\mathrm{lin}(#1)}

% ============================== misc
\newcommand{\E}{\mathbb{E}}
\newcommand{\Prob}{\mathbb{P}}
\newcommand{\indicator}[1]{\mathbb{1}\!\left[#1\right]}
